% GENERATED FILE. Edit canonical lessons, then run npm run generate:books.
% canonical-source-hash: fnv1a64:32bcd6917350679b
% canonical-lessons: JA-C44-shima, JA-C44-nami, JA-C44-tsuchi, JA-C44-tane, JA-C44-ame

\chapter{Island, Wave, Soil, Seed, Rain}
\label{ch:ja-island-wave-soil-seed-rain}
\hlchaptermodality{44}

\begin{chapteropening}
\textbf{By the end of this chapter:} \emph{I can say an island, a wave, soil, a seed and rain.}

Complete the last lesson of chapter 44: i can say an island, a wave, soil, a seed and rain.
\end{chapteropening}

\section[shima]{\ja{しま} (shima) — an island}

\label{lesson:JA-C44-shima}



\begin{quote}
Before the new one: say the Japanese for a cloud; a spider, then the Japanese for the sun.
\end{quote}

\subsection*{You'll want to know: \ja{しま}}
\textbf{\ja{しま}} — \emph{shima} — \textquotedblleft{}an island\textquotedblright{}.

\begin{grammarlens}[title={What you've built}]
One more word for this chapter.
\end{grammarlens}

\subsection*{Your turn}
\begin{itemize}
\raggedright
  \item \emph{Say it:} \emph{shima}
  \item \emph{Say it:} \emph{shima}, once more
  \item \emph{Say it:} say \emph{shima}
  \item \emph{Recall:} say the Japanese for a cloud; a spider, then the Japanese for the sun, then say \emph{shima} again
\end{itemize}

\begin{tcolorbox}[breakable,colback=teal!4,colframe=teal!35!black,title={Before you move on}]
What does \ja{しま} mean? (\textquotedblleft{}An island\textquotedblright{}.) Say it once more.
\end{tcolorbox}
\section[nami]{\ja{なみ} (nami) — a wave}

\label{lesson:JA-C44-nami}



\begin{quote}
Before the new one: say the Japanese for the sun, then the Japanese for an island.
\end{quote}

\subsection*{You'll want to know: \ja{なみ}}
\textbf{\ja{なみ}} — \emph{nami} — \textquotedblleft{}a wave\textquotedblright{}.

\begin{grammarlens}[title={What you've built}]
One more word for this chapter.
\end{grammarlens}

\subsection*{Your turn}
\begin{itemize}
\raggedright
  \item \emph{Say it:} \emph{nami}
  \item \emph{Say it:} \emph{nami}, once more
  \item \emph{Say it:} say \emph{nami}
  \item \emph{Recall:} say the Japanese for the sun, then the Japanese for an island, then say \emph{nami} again
\end{itemize}

\begin{tcolorbox}[breakable,colback=teal!4,colframe=teal!35!black,title={Before you move on}]
What does \ja{なみ} mean? (\textquotedblleft{}A wave\textquotedblright{}.) Say it once more.
\end{tcolorbox}
\section[tsuchi]{\ja{つち} (tsuchi) — soil}

\label{lesson:JA-C44-tsuchi}



\begin{quote}
Before the new one: say the Japanese for an island, then the Japanese for a wave.
\end{quote}

\subsection*{You'll want to know: \ja{つち}}
\textbf{\ja{つち}} — \emph{tsuchi} — \textquotedblleft{}soil\textquotedblright{}.

\begin{grammarlens}[title={What you've built}]
One more word for this chapter.
\end{grammarlens}

\subsection*{Your turn}
\begin{itemize}
\raggedright
  \item \emph{Say it:} \emph{tsuchi}
  \item \emph{Say it:} \emph{tsuchi}, once more
  \item \emph{Say it:} say \emph{tsuchi}
  \item \emph{Recall:} say the Japanese for an island, then the Japanese for a wave, then say \emph{tsuchi} again
\end{itemize}

\begin{tcolorbox}[breakable,colback=teal!4,colframe=teal!35!black,title={Before you move on}]
What does \ja{つち} mean? (\textquotedblleft{}Soil\textquotedblright{}.) Say it once more.
\end{tcolorbox}
\section[tane]{\ja{たね} (tane) — a seed}

\label{lesson:JA-C44-tane}



\begin{quote}
Before the new one: say the Japanese for a wave, then the Japanese for soil.
\end{quote}

\subsection*{You'll want to know: \ja{たね}}
\textbf{\ja{たね}} — \emph{tane} — \textquotedblleft{}a seed\textquotedblright{}.

\begin{grammarlens}[title={What you've built}]
One more word for this chapter.
\end{grammarlens}

\subsection*{Your turn}
\begin{itemize}
\raggedright
  \item \emph{Say it:} \emph{tane}
  \item \emph{Say it:} \emph{tane}, once more
  \item \emph{Say it:} say \emph{tane}
  \item \emph{Recall:} say the Japanese for a wave, then the Japanese for soil, then say \emph{tane} again
\end{itemize}

\begin{tcolorbox}[breakable,colback=teal!4,colframe=teal!35!black,title={Before you move on}]
What does \ja{たね} mean? (\textquotedblleft{}A seed\textquotedblright{}.) Say it once more.
\end{tcolorbox}
\section[ame]{\ja{あめ} (ame) — rain; candy}

\label{lesson:JA-C44-ame}



\begin{quote}
Before the new one: say the Japanese for soil, then the Japanese for a seed.
\end{quote}

\subsection*{You'll want to know: \ja{あめ}}
\textbf{\ja{あめ}} — \emph{ame} — \textquotedblleft{}rain; candy\textquotedblright{}.

\begin{grammarlens}[title={What you've built}]
One more word for this chapter.
\end{grammarlens}

\subsection*{Your turn}
\begin{itemize}
\raggedright
  \item \emph{Say it:} \emph{ame}
  \item \emph{Say it:} \emph{ame}, once more
  \item \emph{Say it:} say \emph{ame}
  \item \emph{Recall:} say the Japanese for soil, then the Japanese for a seed, then say \emph{ame} again
\end{itemize}

\begin{tcolorbox}[breakable,colback=teal!4,colframe=teal!35!black,title={Before you move on}]
What does \ja{あめ} mean? (\textquotedblleft{}Rain; candy\textquotedblright{}.) Say it once more.
\end{tcolorbox}
